% Author: Anachuri Nicolas Daniel
% Template for myself
 
\documentclass[11pt]{article}
\usepackage{indentfirst} 		% First line paragraph indentation
\usepackage{etoolbox} 			% Used for being able to utilise if statements in the language detection.
\setlength{\parskip}{\baselineskip}	% With this line, it is not needed the use of \\ at the end of each paragraph for spacing purposes


% ========================= VARIABLES TO MODIFY =================================

\def\LANGUAGE{EN} %EN, ES		%Language in which the document is going to be written.

\def\reportTile{T. Tecnicas de Comunicacion - Diagnostico }			%Document's title
\def\subject{Trabajo Practico N.1}			%Subject

\def\writter{Anachuri Nicolas Daniel}		%Author(s)


\def\leftUpperHeader{\subject}		%\subject
\def\rightUpperHeader{\leftmark}	%\leftmark for current section


% Do NOT modify anything from here until the beginning of the document.
%============================== DATA PROCESSING ============================
\ifdefstring{\LANGUAGE}{EN}
{

\def\course{1ro del Superior Informatica} %---------
\def\university{Escuela de Minas Dr. Horacio Carrillo} %------------
\def\pageCounterName{PAGE} 		% In the footer will be shown this... 
\def\pageSeparator{OUT OF} 		% ...text in addition of the page number.

}
{
\usepackage[utf8]{inputenc}		%Tildes en caso de no usar arch
\usepackage[spanish, es-tabla]{babel} 	%La opción es-tabla, hace que por defecto las tablas se llamen "Tabla", en vez de "Cuadro".
\def\course{1ro del Superior Informatica} %----------
\def\university{Escuela de Minas Dr. Horacio Carrillo}
\def\pageCounterName{PÁGINA}		%En el pie de página aparece este contenido más el número de página
\def\pageSeparator{DE} 			%DE, OUT OF

}


%===================== USEFUL VARIABLES  =============================
\def\imageSize{0.9}

%===================== PACKAGES TO USE  =============================
\usepackage{amsmath} 			% Allows me the use of matrix in math mode
\setcounter{MaxMatrixCols}{20} 		% Increases the maximum number of columns in matrix from 10 to 20.
\usepackage{pdfpages} 			% Attach pdfs using \includepdf[pages=initial-final]{path.pdf}
\usepackage{lastpage}			% Used to reference the last page and include it in the footer.
\usepackage{xcolor}			% Change color of text and so on. \textcolor{color}{text} is the one I use the most.
\usepackage{placeins} 			% Allows the use of the command \FloatBarrier
\usepackage{setspace}			% Allows the use of the spacing environment to change the space between lines.
\usepackage{graphicx}			% Figures addition
\usepackage{wrapfig}
\usepackage{geometry}			% Changes the geometry of the pages
\usepackage[small, bf]{caption}		% Decreases the size of captions and turns them bold text.
\usepackage{subcaption}			% Include subfigures.
\usepackage{hyperref}			% Allows clickable references.
\hypersetup{colorlinks=true, allcolors=black} 	% Colors of links
\usepackage{pdflscape}			% Allows the use of the environment lscape for landscape pages

\usepackage{fancyhdr}			% Modify header and footer
\usepackage{bm}				% Allows to use bold text in math mode with \bm
\usepackage[makeroom]{cancel}		% Allows the use of \cancelto{}{} to cross out equations
\usepackage{titlesec}			% Change title format


%========== CHANGE TITLE FORMAT  ======================
%\titleformat{\section}
%{\bfseries}	% format
%{\thesection}	% label
%{0.3cm}		% separation between label and body
%{}		% code preceding title body
%[]		% code following title body


%========== HEADER AND FOOTER CONFIGURATION ======================
\pagestyle{fancy}
% HEADER[EVEN PAGES]{ODD PAGES}
% Left header
\lhead[
	\scriptsize{\MakeUppercase{\leftUpperHeader}} % Even pages
]
{
	\scriptsize{\MakeUppercase{\leftUpperHeader}} % Odd pages
}

% Central header
\chead[]{}

% Right header
\rhead[
\scriptsize{\rightUpperHeader} % Even pages
]
{
\scriptsize{\rightUpperHeader} % Odd pages
}

\renewcommand{\headrulewidth}{0.8pt}	% Width of the header rule

% FOOTER[EVEN PAGES]{ODD PAGES}
% Left footer
\lfoot[]{}

% Central footer
\cfoot[
\tiny{\pageCounterName\space\thepage\space \pageSeparator\space\pageref{LastPage}} % Even pages
]
{
\tiny{\pageCounterName\space\thepage\space \pageSeparator\space\pageref{LastPage}} % Odd pages
}

% Right footer
\rfoot[]{}
\renewcommand{\footrulewidth}{0.8pt} %Width of the footer rule


%================================== USER OWN FUNCTIONS ===============================================
\newcommand{\PDEA}[1]{\cdot 10^{#1}} % Función para escribir más rápidamente multiplicaciones por potencias de 10. PDEA = Por Diez Elevado A

\usepackage{xargs}	%Permite manejar argumentos opcionales en los comandos creados por el usuario

%Incluir imágenes, la sintaxis es la siguiente:
%\IncludeImage{ruta}[escalado][pie figura][etiqueta], los corchetes son opcionales
\newcommandx*{\IncludeImage}[4][2=1, 3=, 4=]{
	%#1 es la ruta a la imagen
	%#2 es el escalado
	%#3 es el pie de figura
	%#4 es la etiqueta
	
	\begin{figure}[h!]
		\centering
		\includegraphics[width=#2\linewidth]{#1}
		%Si se especifica pie de figura...
		\ifblank{#3}{}{
			\caption{#3}	
		}
		%Si se especifica etiqueta...
		\ifblank{#4}{}{
			\label{#4}	
		}
	\end{figure}
	\FloatBarrier
}



%================================ DOCUMENT BEGINNING  =============================================
\begin{document}
%***************** TITLE PAGE (DO NOT MODIFY) ******************************
\begin{titlepage}
	\newgeometry{top=2cm, bottom=3.5cm}

	%Logos Escuela y universidad
	\def\logoSize{0.2}
	\begin{figure}
	\includegraphics[width=\logoSize\linewidth]{/home/nico/Descargas/minas.jpeg}
	\hfill
	\includegraphics[width=\logoSize\linewidth]{/home/nico/Descargas/unju.jpeg}
	\end{figure}


	%Espacio entre logos y título
	\vspace*{1.75cm}

	%Título con interlineado aumentado
	\begin{spacing}{2}
	\centering{
	\Huge{
		\textbf{\reportTile}
	}
	}
	\end{spacing}
	
	%Línea horizontal de la portada
	\hrule
	
	%Se baja al final de la hoja
	\vfill
	\begin{flushright}
	\LARGE{\writter}

	\vspace{2cm}
	
	\Large
	\subject\\
	\course\\
	\university
	
	\vspace{1cm}
	
	\today
	\end{flushright}
	
\end{titlepage}
	
\newgeometry{top=3cm, bottom=3cm}

\tableofcontents
%\cleardoublepage
\clearpage


\section{Propone una solución para las siguientes situaciones enigmáticas}

\begin{itemize}
    
  \item Si hay un accidente en la frontera entre México y Estados Unidos, ¿dónde se entierran los sobrevivientes?
  
  ¿Quien en su sano juicio enterraría a personas vivas?
  
  \item ¿Es legal que un hombre contraiga matrimonio con la hermana de su viuda?
    
   Es posible teniendo en cuenta que de alguna forma inexplicable el hombre pueda comunicarse con la hermana desde "el otro lado" y sea legal casarse con un espectro. De otra forma es totalmente imposible que alguien muerto pueda contraer matrimonio otra vez.  
  
  \item ¿Cuánta tierra puedes encontrar en un hoyo de 5 metros de profundidad y 2 de ancho?  
  
  Dada la ubicación del hoyo es posible encontrar partículas de tierra. 
  
  \item ¿Cuántas manzanas crecen en un árbol?
    
  Según mi conocimiento todas las manzanas crecen en el árbol. 
    
  \item  Resuelve este acertijo: Cual es la palabra de 4 letras, que tiene 3 letras, aunque se escribe con 6, raramente consta de 9, mientras tiene 8, y nunca se escribe con 5.   
    
   En el próximo Trabajo Práctico responderé la pregunta, ahora mismo tengo pistas, pero todavía se me es imposible. 
  
\end{itemize}



\section{Explosión de Sabores}

En el momento de querer expresar mis emociones las palabras simplemente no son capaces de salir de mi garganta, mi mente de niebla y mi cuerpo se congela, pero, cuando estoy en presencia de este producto  traído por los mismísimos ángeles, el cual se puede mostrar de dientes formas: negro, blanco, con leche, etc. sin embargo ninguno de ellos prevalece sobre otro, aunque diga esto, lamentablemente tengo un favorito el Chocolate negro, porque? Su sabor amargo y dulce mezclándose al mismo tiempo en mi paladar no tiene precio ni explicación, una explosión de frescura y tranquilidad se expande por tu mi ser.

Para expresar de una forma entendible para otro ser vivo este tipo de sensaciones debo recurrir a palabras, palabras... ¿por donde debería empezar? Uno solo pensaría que es el sabor, pero también entra en juego la forma, el aroma, el tacto, la envoltura. Pero el momento en el que se derrite en tu boca es único e inolvidable, las palabras se agotan para expresarlo por este motivo recurriré a una imagen que lo describe perfectamente:

\clearpage

\begin{wrapfigure}{l}{0.6\linewidth}
  

  \centering
  \includegraphics[width=0.5\linewidth]{/home/nico/me/Comunicacion/TP-1-AnachuriNicolas/feel.jpg}
  \caption{Explosion de Sentimientos}
  \label{fig:etiqueta}

\end{wrapfigure}

Una combinación de felicidad, placer, relajación, alivio, Tristeza porque sé que en algún momento estos sentimientos se acabaran y envidia por aquellos que pueden comer cuando lo deseen. 





\end{document}

